% Synthetic theory paper for proof_index.py. Has a clean dependency chain
% plus one deliberately undefined ref.
\documentclass{article}
\usepackage{amsmath, amsthm}
\newtheorem{theorem}{Theorem}
\newtheorem{lemma}{Lemma}
\newtheorem{proposition}{Proposition}
\newtheorem{assumption}{Assumption}

\begin{document}

\begin{assumption}\label{ass:subgauss}
The noise is sub-Gaussian with parameter sigma.
\end{assumption}

\begin{lemma}\label{lem:concentration}
Under Assumption~\ref{ass:subgauss}, the empirical mean concentrates.
\end{lemma}
\begin{proof}
By Assumption~\ref{ass:subgauss} and a Chernoff bound. See also \ref{ass:subgauss}.
\end{proof}

\begin{lemma}\label{lem:bias}
The bias term is bounded by a smoothing argument.
\end{lemma}
\begin{proof}
Direct calculation. No dependencies.
\end{proof}

\begin{theorem}\label{thm:rate}
The estimator achieves the minimax rate.
\end{theorem}
\begin{proof}
Combine Lemma~\ref{lem:concentration} and Lemma~\ref{lem:bias}.
We also use Assumption~\ref{ass:subgauss}.
\end{proof}

\begin{proposition}\label{prop:extra}
A corollary-style extension.
\end{proposition}
\begin{proof}
Follows from Theorem~\ref{thm:rate}. Also cites a nonexistent label
\ref{thm:does_not_exist} which must be flagged.
\end{proof}

% An unlabeled lemma (should WARN).
\begin{lemma}
An auxiliary bound stated without a label.
\end{lemma}

\end{document}
